% META: {"id": "1NSI-TB-ME-004", "chapitre": "1NSI-TYPES-BASE", "type_objet": "methode", "capacites": ["P-BASE-05"], "capacites_codes": ["C5"], "methodes": ["M4"], "mode_creation": "ex_nihilo", "sources_inspiration": [], "status": "needs_review"}
% BEGIN-VERIFY
% assert ord('A') == 65
% assert chr(233) == 'é'
% assert len("Bonjour a tous !") == 16
% assert len("Bonjour a tous !".encode("ascii")) == 16
% assert len("é".encode("utf-8")) == 2
% assert len("é".encode("iso-8859-1")) == 1
% assert "Bonjour".encode("ascii") == "Bonjour".encode("utf-8")
% try:
%     "é".encode("ascii")
%     ascii_accepte_accent = True
% except UnicodeEncodeError:
%     ascii_accepte_accent = False
% assert ascii_accepte_accent is False
% END-VERIFY
\begin{fichemethode}{M4}{Choisir un encodage de texte et justifier ce choix}

\textbf{Quand l'utiliser ?} Quand un énoncé demande quel encodage convient à un texte, ou
pourquoi des caractères s'affichent de travers.

\textbf{Pas à pas :}
\begin{enumerate}
  \item \textbf{Regarde le texte, caractère par caractère.} La question n'est pas « quel
        encodage est le meilleur » mais « lesquels peuvent représenter \emph{ces}
        caractères-là ».
  \item \textbf{Situe ASCII} : $128$ caractères, l'alphabet latin non accentué, les chiffres et
        la ponctuation courante. Un seul octet par caractère.
  \item \textbf{Situe les autres} : ISO-8859-1 ajoute les caractères accentués d'Europe
        occidentale, toujours sur un octet. UTF-8 représente \emph{tous} les caractères, sur
        un à quatre octets selon le caractère.
  \item \textbf{Distingue « peut représenter » de « est seul à pouvoir ».} Plusieurs encodages
        conviennent souvent : les $128$ premiers caractères d'UTF-8 sont exactement ceux
        d'ASCII, et un texte sans accent s'encode identiquement dans les trois.
  \item \textbf{Compte les octets, pas les caractères}, quand la question porte sur la taille.
\end{enumerate}

\textbf{Exemple rédigé.} « Bonjour à tous ! » contient un « à » : ASCII ne suffit pas.
ISO-8859-1 convient, UTF-8 aussi. En UTF-8 le « à » occupe \textbf{deux} octets, contre un
seul en ISO-8859-1 : même texte, tailles différentes.
Sans accent, « Bonjour » s'encode à l'identique dans les trois.

\textbf{Erreur fréquente.} Retenir qu'UTF-8 encode tout — ce qui est vrai — et en conclure
qu'il est le seul à convenir. Ou oublier qu'ASCII et UTF-8 conviennent aussi dès que le texte
n'a pas d'accent.

\end{fichemethode}
